\chapter{Bijektionen auf Funktionsmengen}
\label{b08guide:ch:functionsets}

Wie in den Grundbegriffen bezeichnet $B^A$ die Menge aller
Funktionen von $A$ nach $B$ und $\operatorname{Bij}(A,B)$ die
Menge der Bijektionen zwischen diesen angegebenen Mengen.
Zusätzliche Strukturen auf $A$ und $B$ werden nicht vorausgesetzt.

\section{Vor- und Nachkomposition}

\begin{theorem}[Vor- und Nachkomposition mit einer Bijektion]
\label{b08guide:thm:pre-post-composition}
Seien $u\colon A\bij A'$ und $v\colon B\bij B'$ Bijektionen
und $C$ eine beliebige Menge. Dann sind
\[
 \begin{aligned}
 \mathcal C_u\colon C^{A'}&\bij C^A,
 &\mathcal C_u(h)&=h\circ u,\\
 \mathcal D_v\colon B^C&\bij (B')^C,
 &\mathcal D_v(h)&=v\circ h
 \end{aligned}
\]
Bijektionen. Ihre Umkehrungen lauten
\[
 (\mathcal C_u)^{-1}(k)=k\circ u^{-1},
 \qquad (\mathcal D_v)^{-1}(k)=v^{-1}\circ k.
\]
Alle beteiligten Mengen dürfen leer sein, soweit die angegebenen
Bijektionen existieren.
\end{theorem}

\begin{theorem}[Gleichzeitiger Wechsel von Definitions- und Zielmenge]
\label{b08guide:thm:function-transport}
Für Bijektionen $u\colon A\bij A'$ und $v\colon B\bij B'$ ist
\[
 \mathcal T_{u,v}\colon B^A\bij (B')^{A'},
 \qquad h\mapsto v\circ h\circ u^{-1},
\]
eine Bijektion. Die Umkehrung ist
\[
 \mathcal T_{u,v}^{-1}(k)=v^{-1}\circ k\circ u.
\]
Insbesondere entspricht jeder Funktion $h\colon A\to B$
genau eine Funktion $k\colon A'\to B'$ mit
\[
 k\circ u=v\circ h.
\]
\end{theorem}

\section{Mengen von Bijektionen}

\begin{theorem}[Übertragung ganzer Bijektionsmengen]
\label{b08guide:thm:bijection-transport}
Für Bijektionen $u\colon A\bij A'$ und $v\colon B\bij B'$
beschränkt sich die Funktion aus
Satz~\ref{b08guide:thm:function-transport} zu einer Bijektion
\[
 \mathcal T_{u,v}\big|_{\operatorname{Bij}(A,B)}
 \colon\operatorname{Bij}(A,B)\bij\operatorname{Bij}(A',B'),
 \qquad h\mapsto v\circ h\circ u^{-1}.
\]
Ihre Umkehrung ist wiederum $k\mapsto v^{-1}\circ k\circ u$.
Insbesondere gilt
\[
 \operatorname{Bij}(A,B)\ne\varnothing
 \quad\Longleftrightarrow\quad
 \operatorname{Bij}(A',B')\ne\varnothing.
\]
Es wird keine Bijektion zwischen $A$ und $B$ vorausgesetzt;
beide Bijektionsmengen können leer sein.
\end{theorem}

\begin{theorem}[Umkehrzuordnung zwischen Bijektionsmengen]
\label{b08guide:thm:inverse-assignment}
Für beliebige Mengen $A,B$ ist
\[
 \mathcal I_{A,B}\colon\operatorname{Bij}(A,B)
                  \bij\operatorname{Bij}(B,A),
 \qquad f\mapsto f^{-1},
\]
eine Bijektion. Es gilt
\[
 \mathcal I_{A,B}^{-1}=\mathcal I_{B,A}.
\]
Falls keine Bijektion zwischen $A$ und $B$ existiert,
ist $\mathcal I_{A,B}$ die leere Bijektion. Für $A=B$ ist
die zweimalige Anwendung der Umkehrzuordnung die Identität
auf $\operatorname{Bij}(A,A)$.
\end{theorem}

\section{Ein Test mit genau zwei Werten}

\begin{theorem}[Charakterisierung durch eine zweielementige Wertemenge]
\label{b08guide:thm:two-value-test}
Sei $\mathbf 2=\{0,1\}$ eine feste Menge mit $0\ne1$.
Für eine beliebige Funktion $f\colon A\to B$ sind folgende
Aussagen äquivalent:
\begin{enumerate}
 \item $f\colon A\bij B$ ist bijektiv.
 \item Die Präkomposition
 \[
  \mathcal C_f\colon\mathbf 2^B\to\mathbf 2^A,
  \qquad g\mapsto g\circ f,
 \]
 ist bijektiv.
 \item Für jede Funktion $h\colon A\to\mathbf 2$ gibt es
 genau eine Funktion $g\colon B\to\mathbf 2$ mit $h=g\circ f$.
\end{enumerate}
Der Test gilt einschließlich aller leeren Randfälle und ohne
Auswahlaxiom. Die Elemente $0$ und $1$ dienen allein als zwei
verschiedene Werte; eine Rechen- oder Ordnungsstruktur wird
nicht verwendet.
\end{theorem}
